% 10_boot_security.tex -- secure boot chain, signature check, TrustZone, security blocks, JTAG debug, FPGA configuration.
% Build: tectonic 10_boot_security.tex
\documentclass[tikz,border=10pt,12pt]{standalone}
% soc_style.tex -- shared look for the SoC figures (same palette as the slide decks and the
% core register placement figure). Every figure does % soc_style.tex -- shared look for the SoC figures (same palette as the slide decks and the
% core register placement figure). Every figure does % soc_style.tex -- shared look for the SoC figures (same palette as the slide decks and the
% core register placement figure). Every figure does \input{soc_style}.
\usepackage[T1]{fontenc}
\usepackage[scaled=0.95]{helvet}
\renewcommand{\familydefault}{\sfdefault}
\usepackage{booktabs}
\usepackage{amsmath}
\usepackage{amssymb}
\usetikzlibrary{arrows.meta,calc,positioning,fit,shapes.geometric,shapes.misc}

% ---------------- palette ----------------
\definecolor{pcb}{HTML}{0B2E33}
\definecolor{pcbmid}{HTML}{1F5C61}
\definecolor{tint}{HTML}{E6F0EF}
\definecolor{copper}{HTML}{B5651D}
\definecolor{coppert}{HTML}{F6E8DC}
\definecolor{ink}{HTML}{1F2A2E}
\definecolor{muted}{HTML}{5B6B70}
\definecolor{hardf}{HTML}{DCE6F2}
\definecolor{harde}{HTML}{2F4A7A}
\definecolor{good}{HTML}{2E7D32}
\definecolor{bad}{HTML}{B3261E}
% one colour per kind of block: border (d) / fill (f)
\definecolor{cpud}{HTML}{2B5C9E} \definecolor{cpuf}{HTML}{D6E4F5}
\definecolor{gpud}{HTML}{2E7D32} \definecolor{gpuf}{HTML}{D5EDD0}
\definecolor{npud}{HTML}{0B2E33} \definecolor{npuf}{HTML}{BFE0DB}
\definecolor{memd}{HTML}{5E3A8C} \definecolor{memf}{HTML}{E4DAF3}
\definecolor{busd}{HTML}{B5651D} \definecolor{busf}{HTML}{F6E8DC}
\definecolor{prd}{HTML}{546E7A}  \definecolor{prf}{HTML}{E3EAED}
\definecolor{secd}{HTML}{B3261E} \definecolor{secf}{HTML}{F7D9D6}
\definecolor{pwrd}{HTML}{8A6800} \definecolor{pwrf}{HTML}{FFF0B8}
\definecolor{medd}{HTML}{8E2463} \definecolor{medf}{HTML}{F5D6E6}
\definecolor{skd}{HTML}{2E7D32}  \definecolor{skf}{HTML}{BFE3B4}   % skew buffers, as in the register figure

\tikzset{
  blk/.style 2 args={draw=#1, fill=#2, rounded corners=4pt, thick, align=center,
                     font=\scriptsize, inner sep=3pt},
  cpu/.style={blk={cpud}{cpuf}},  gpu/.style={blk={gpud}{gpuf}},
  npu/.style={blk={npud}{npuf}},  mem/.style={blk={memd}{memf}},
  bus/.style={blk={busd}{busf}},  per/.style={blk={prd}{prf}},
  sec/.style={blk={secd}{secf}},  pwr/.style={blk={pwrd}{pwrf}},
  med/.style={blk={medd}{medf}},
  comb/.style ={draw=copper, fill=coppert, rounded corners=2pt, thick, align=center,
                font=\scriptsize, inner sep=2pt},                       % combinational logic
  hard/.style ={draw=harde, fill=hardf, thick, dashed, align=center,
                font=\scriptsize, inner sep=2pt},                       % hard macro / off-chip
  sram/.style ={draw=memd, fill=memf, thick, align=center, font=\scriptsize,
                inner sep=2pt},                                         % SRAM array, not flip-flops
  ff/.style   ={draw=#1, fill=#1!30, line width=0.35pt},                % one flip-flop square
  ffgone/.style={draw=muted, fill=white, densely dashed, line width=0.35pt},
  panel/.style={draw=muted, rounded corners=6pt, line width=0.9pt},
  ptitle/.style={anchor=north west, font=\normalsize\bfseries, text=ink, inner sep=0pt},
  note/.style ={font=\scriptsize, text=ink, align=left, inner sep=0pt},
  lab/.style  ={font=\tiny, text=muted, inner sep=1pt},
  axi/.style  ={{Stealth[length=2.6mm]}-{Stealth[length=2.6mm]}, line width=2.2pt, copper},
  axi1/.style ={-{Stealth[length=2.6mm]}, line width=2.2pt, copper},
  apb/.style  ={{Stealth[length=1.8mm]}-{Stealth[length=1.8mm]}, line width=1pt, copper},
  apb1/.style ={-{Stealth[length=1.8mm]}, line width=1pt, copper},
  dat/.style  ={-{Stealth[length=1.8mm]}, thick, pcbmid},
  ctl/.style  ={-{Stealth[length=1.8mm]}, densely dashed, thick, muted},
  irq/.style  ={-{Stealth[length=2mm]}, densely dotted, line width=1.3pt, bad},
  num/.style  ={circle, fill=copper, text=white, font=\scriptsize\bfseries, inner sep=1pt,
                minimum size=0.46cm},
  ours/.style ={draw=bad, line width=1.2pt, densely dashed, rounded corners=4pt},
}

% ---------------- helpers ----------------
\newcommand{\rn}[1]{\texttt{#1}}
% \ffrow{x}{y}{n}{cell}{style}: n flip-flops in a row, bottom-left corner at (x,y)
\newcommand{\ffrow}[5]{\foreach \q in {1,...,#3}{\draw[#5] ({#1+(\q-1)*#4},{#2}) rectangle ++(#4,#4);}}
% \ffcol{x}{y}{n}{cell}{style}: n flip-flops in a column, top-left corner at (x,y)
\newcommand{\ffcol}[5]{\foreach \q in {1,...,#3}{\draw[#5] ({#1},{#2-\q*#4}) rectangle ++(#4,#4);}}
% \ffgrid{x}{ytop}{n}{per row}{cell}{style}: n flip-flops, filled row by row from the top-left
\newcommand{\ffgrid}[6]{\pgfmathtruncatemacro{\gl}{#3-1}%
  \foreach \b in {0,...,\gl}{\pgfmathtruncatemacro{\gc}{mod(\b,#4)}\pgfmathtruncatemacro{\gr}{div(\b,#4)}%
    \draw[#6] ({#1+\gc*#5},{#2-(\gr+1)*#5}) rectangle ++(#5,#5);}}

% every figure: same page size, title, footer
\newcommand{\canvas}{\useasboundingbox (-0.3,-0.2) rectangle (38.6,30.4);}
\newcommand{\figtitle}[3]{%
  \node[num, minimum size=0.9cm, font=\Large\bfseries] at (0.45,29.55) {#1};
  \node[anchor=west, font=\LARGE\bfseries, text=pcb] at (1.15,29.55) {#2};
  \node[anchor=west, font=\small, text=ink] at (0,28.7) {#3};}
\newcommand{\figfoot}[1]{\node[anchor=west, font=\scriptsize, text=muted] at (0,0.25) {#1};}
% \panel{x0}{y0}{x1}{y1}{title}
\newcommand{\panel}[5]{\draw[panel] (#1,#2) rectangle (#3,#4);
  \node[ptitle] at ({#1+0.3},{#4-0.25}) {#5};}

% ---------------- timing diagrams: 1 cycle = \cw cm, low at y, high at y+0.45 ----------------
\newcommand{\cw}{1.0}
\newcommand{\wclk}[3]{% x0 y ncycles
  \foreach \k in {1,...,#3}{\draw[pcb, thick] ({#1+(\k-1)*\cw},#2) -- ++(0,0.45) -- ++(0.5*\cw,0)
     -- ++(0,-0.45) -- ++(0.5*\cw,0);}}
\newcommand{\wsig}[3]{% x0 y {level per cycle, e.g. 0,1,1,0}
  \gdef\wprev{x}%
  \foreach \v [count=\k from 0] in {#3}{%
    \draw[pcb, thick] ({#1+\k*\cw},{#2+0.45*\v}) -- ++(\cw,0);
    \if x\wprev\else\ifnum\wprev=\v\else\draw[pcb, thick] ({#1+\k*\cw},#2) -- ++(0,0.45);\fi\fi
    \global\let\wprev\v}}
\newcommand{\wbus}[5]{% x0 y c0 c1 text: one value held from cycle c0 to cycle c1
  \draw[pcbmid, thick, fill=white] ({#1+(#3)*\cw},{#2+0.225}) -- ({#1+(#3)*\cw+0.08},{#2+0.45})
    -- ({#1+(#4)*\cw-0.08},{#2+0.45}) -- ({#1+(#4)*\cw},{#2+0.225}) -- ({#1+(#4)*\cw-0.08},#2)
    -- ({#1+(#3)*\cw+0.08},#2) -- cycle;
  \node[font=\tiny\ttfamily] at ({#1+((#3)+(#4))*\cw/2},{#2+0.225}) {#5};}
\newcommand{\widle}[4]{% x0 y c0 c1: bus idle (no value)
  \draw[muted, thick] ({#1+(#3)*\cw},{#2+0.225}) -- ({#1+(#4)*\cw},{#2+0.225});}
\newcommand{\wname}[3]{\node[anchor=east, font=\scriptsize\ttfamily] at ({#1-0.15},{#2+0.225}) {#3};}
.
\usepackage[T1]{fontenc}
\usepackage[scaled=0.95]{helvet}
\renewcommand{\familydefault}{\sfdefault}
\usepackage{booktabs}
\usepackage{amsmath}
\usepackage{amssymb}
\usetikzlibrary{arrows.meta,calc,positioning,fit,shapes.geometric,shapes.misc}

% ---------------- palette ----------------
\definecolor{pcb}{HTML}{0B2E33}
\definecolor{pcbmid}{HTML}{1F5C61}
\definecolor{tint}{HTML}{E6F0EF}
\definecolor{copper}{HTML}{B5651D}
\definecolor{coppert}{HTML}{F6E8DC}
\definecolor{ink}{HTML}{1F2A2E}
\definecolor{muted}{HTML}{5B6B70}
\definecolor{hardf}{HTML}{DCE6F2}
\definecolor{harde}{HTML}{2F4A7A}
\definecolor{good}{HTML}{2E7D32}
\definecolor{bad}{HTML}{B3261E}
% one colour per kind of block: border (d) / fill (f)
\definecolor{cpud}{HTML}{2B5C9E} \definecolor{cpuf}{HTML}{D6E4F5}
\definecolor{gpud}{HTML}{2E7D32} \definecolor{gpuf}{HTML}{D5EDD0}
\definecolor{npud}{HTML}{0B2E33} \definecolor{npuf}{HTML}{BFE0DB}
\definecolor{memd}{HTML}{5E3A8C} \definecolor{memf}{HTML}{E4DAF3}
\definecolor{busd}{HTML}{B5651D} \definecolor{busf}{HTML}{F6E8DC}
\definecolor{prd}{HTML}{546E7A}  \definecolor{prf}{HTML}{E3EAED}
\definecolor{secd}{HTML}{B3261E} \definecolor{secf}{HTML}{F7D9D6}
\definecolor{pwrd}{HTML}{8A6800} \definecolor{pwrf}{HTML}{FFF0B8}
\definecolor{medd}{HTML}{8E2463} \definecolor{medf}{HTML}{F5D6E6}
\definecolor{skd}{HTML}{2E7D32}  \definecolor{skf}{HTML}{BFE3B4}   % skew buffers, as in the register figure

\tikzset{
  blk/.style 2 args={draw=#1, fill=#2, rounded corners=4pt, thick, align=center,
                     font=\scriptsize, inner sep=3pt},
  cpu/.style={blk={cpud}{cpuf}},  gpu/.style={blk={gpud}{gpuf}},
  npu/.style={blk={npud}{npuf}},  mem/.style={blk={memd}{memf}},
  bus/.style={blk={busd}{busf}},  per/.style={blk={prd}{prf}},
  sec/.style={blk={secd}{secf}},  pwr/.style={blk={pwrd}{pwrf}},
  med/.style={blk={medd}{medf}},
  comb/.style ={draw=copper, fill=coppert, rounded corners=2pt, thick, align=center,
                font=\scriptsize, inner sep=2pt},                       % combinational logic
  hard/.style ={draw=harde, fill=hardf, thick, dashed, align=center,
                font=\scriptsize, inner sep=2pt},                       % hard macro / off-chip
  sram/.style ={draw=memd, fill=memf, thick, align=center, font=\scriptsize,
                inner sep=2pt},                                         % SRAM array, not flip-flops
  ff/.style   ={draw=#1, fill=#1!30, line width=0.35pt},                % one flip-flop square
  ffgone/.style={draw=muted, fill=white, densely dashed, line width=0.35pt},
  panel/.style={draw=muted, rounded corners=6pt, line width=0.9pt},
  ptitle/.style={anchor=north west, font=\normalsize\bfseries, text=ink, inner sep=0pt},
  note/.style ={font=\scriptsize, text=ink, align=left, inner sep=0pt},
  lab/.style  ={font=\tiny, text=muted, inner sep=1pt},
  axi/.style  ={{Stealth[length=2.6mm]}-{Stealth[length=2.6mm]}, line width=2.2pt, copper},
  axi1/.style ={-{Stealth[length=2.6mm]}, line width=2.2pt, copper},
  apb/.style  ={{Stealth[length=1.8mm]}-{Stealth[length=1.8mm]}, line width=1pt, copper},
  apb1/.style ={-{Stealth[length=1.8mm]}, line width=1pt, copper},
  dat/.style  ={-{Stealth[length=1.8mm]}, thick, pcbmid},
  ctl/.style  ={-{Stealth[length=1.8mm]}, densely dashed, thick, muted},
  irq/.style  ={-{Stealth[length=2mm]}, densely dotted, line width=1.3pt, bad},
  num/.style  ={circle, fill=copper, text=white, font=\scriptsize\bfseries, inner sep=1pt,
                minimum size=0.46cm},
  ours/.style ={draw=bad, line width=1.2pt, densely dashed, rounded corners=4pt},
}

% ---------------- helpers ----------------
\newcommand{\rn}[1]{\texttt{#1}}
% \ffrow{x}{y}{n}{cell}{style}: n flip-flops in a row, bottom-left corner at (x,y)
\newcommand{\ffrow}[5]{\foreach \q in {1,...,#3}{\draw[#5] ({#1+(\q-1)*#4},{#2}) rectangle ++(#4,#4);}}
% \ffcol{x}{y}{n}{cell}{style}: n flip-flops in a column, top-left corner at (x,y)
\newcommand{\ffcol}[5]{\foreach \q in {1,...,#3}{\draw[#5] ({#1},{#2-\q*#4}) rectangle ++(#4,#4);}}
% \ffgrid{x}{ytop}{n}{per row}{cell}{style}: n flip-flops, filled row by row from the top-left
\newcommand{\ffgrid}[6]{\pgfmathtruncatemacro{\gl}{#3-1}%
  \foreach \b in {0,...,\gl}{\pgfmathtruncatemacro{\gc}{mod(\b,#4)}\pgfmathtruncatemacro{\gr}{div(\b,#4)}%
    \draw[#6] ({#1+\gc*#5},{#2-(\gr+1)*#5}) rectangle ++(#5,#5);}}

% every figure: same page size, title, footer
\newcommand{\canvas}{\useasboundingbox (-0.3,-0.2) rectangle (38.6,30.4);}
\newcommand{\figtitle}[3]{%
  \node[num, minimum size=0.9cm, font=\Large\bfseries] at (0.45,29.55) {#1};
  \node[anchor=west, font=\LARGE\bfseries, text=pcb] at (1.15,29.55) {#2};
  \node[anchor=west, font=\small, text=ink] at (0,28.7) {#3};}
\newcommand{\figfoot}[1]{\node[anchor=west, font=\scriptsize, text=muted] at (0,0.25) {#1};}
% \panel{x0}{y0}{x1}{y1}{title}
\newcommand{\panel}[5]{\draw[panel] (#1,#2) rectangle (#3,#4);
  \node[ptitle] at ({#1+0.3},{#4-0.25}) {#5};}

% ---------------- timing diagrams: 1 cycle = \cw cm, low at y, high at y+0.45 ----------------
\newcommand{\cw}{1.0}
\newcommand{\wclk}[3]{% x0 y ncycles
  \foreach \k in {1,...,#3}{\draw[pcb, thick] ({#1+(\k-1)*\cw},#2) -- ++(0,0.45) -- ++(0.5*\cw,0)
     -- ++(0,-0.45) -- ++(0.5*\cw,0);}}
\newcommand{\wsig}[3]{% x0 y {level per cycle, e.g. 0,1,1,0}
  \gdef\wprev{x}%
  \foreach \v [count=\k from 0] in {#3}{%
    \draw[pcb, thick] ({#1+\k*\cw},{#2+0.45*\v}) -- ++(\cw,0);
    \if x\wprev\else\ifnum\wprev=\v\else\draw[pcb, thick] ({#1+\k*\cw},#2) -- ++(0,0.45);\fi\fi
    \global\let\wprev\v}}
\newcommand{\wbus}[5]{% x0 y c0 c1 text: one value held from cycle c0 to cycle c1
  \draw[pcbmid, thick, fill=white] ({#1+(#3)*\cw},{#2+0.225}) -- ({#1+(#3)*\cw+0.08},{#2+0.45})
    -- ({#1+(#4)*\cw-0.08},{#2+0.45}) -- ({#1+(#4)*\cw},{#2+0.225}) -- ({#1+(#4)*\cw-0.08},#2)
    -- ({#1+(#3)*\cw+0.08},#2) -- cycle;
  \node[font=\tiny\ttfamily] at ({#1+((#3)+(#4))*\cw/2},{#2+0.225}) {#5};}
\newcommand{\widle}[4]{% x0 y c0 c1: bus idle (no value)
  \draw[muted, thick] ({#1+(#3)*\cw},{#2+0.225}) -- ({#1+(#4)*\cw},{#2+0.225});}
\newcommand{\wname}[3]{\node[anchor=east, font=\scriptsize\ttfamily] at ({#1-0.15},{#2+0.225}) {#3};}
.
\usepackage[T1]{fontenc}
\usepackage[scaled=0.95]{helvet}
\renewcommand{\familydefault}{\sfdefault}
\usepackage{booktabs}
\usepackage{amsmath}
\usepackage{amssymb}
\usetikzlibrary{arrows.meta,calc,positioning,fit,shapes.geometric,shapes.misc}

% ---------------- palette ----------------
\definecolor{pcb}{HTML}{0B2E33}
\definecolor{pcbmid}{HTML}{1F5C61}
\definecolor{tint}{HTML}{E6F0EF}
\definecolor{copper}{HTML}{B5651D}
\definecolor{coppert}{HTML}{F6E8DC}
\definecolor{ink}{HTML}{1F2A2E}
\definecolor{muted}{HTML}{5B6B70}
\definecolor{hardf}{HTML}{DCE6F2}
\definecolor{harde}{HTML}{2F4A7A}
\definecolor{good}{HTML}{2E7D32}
\definecolor{bad}{HTML}{B3261E}
% one colour per kind of block: border (d) / fill (f)
\definecolor{cpud}{HTML}{2B5C9E} \definecolor{cpuf}{HTML}{D6E4F5}
\definecolor{gpud}{HTML}{2E7D32} \definecolor{gpuf}{HTML}{D5EDD0}
\definecolor{npud}{HTML}{0B2E33} \definecolor{npuf}{HTML}{BFE0DB}
\definecolor{memd}{HTML}{5E3A8C} \definecolor{memf}{HTML}{E4DAF3}
\definecolor{busd}{HTML}{B5651D} \definecolor{busf}{HTML}{F6E8DC}
\definecolor{prd}{HTML}{546E7A}  \definecolor{prf}{HTML}{E3EAED}
\definecolor{secd}{HTML}{B3261E} \definecolor{secf}{HTML}{F7D9D6}
\definecolor{pwrd}{HTML}{8A6800} \definecolor{pwrf}{HTML}{FFF0B8}
\definecolor{medd}{HTML}{8E2463} \definecolor{medf}{HTML}{F5D6E6}
\definecolor{skd}{HTML}{2E7D32}  \definecolor{skf}{HTML}{BFE3B4}   % skew buffers, as in the register figure

\tikzset{
  blk/.style 2 args={draw=#1, fill=#2, rounded corners=4pt, thick, align=center,
                     font=\scriptsize, inner sep=3pt},
  cpu/.style={blk={cpud}{cpuf}},  gpu/.style={blk={gpud}{gpuf}},
  npu/.style={blk={npud}{npuf}},  mem/.style={blk={memd}{memf}},
  bus/.style={blk={busd}{busf}},  per/.style={blk={prd}{prf}},
  sec/.style={blk={secd}{secf}},  pwr/.style={blk={pwrd}{pwrf}},
  med/.style={blk={medd}{medf}},
  comb/.style ={draw=copper, fill=coppert, rounded corners=2pt, thick, align=center,
                font=\scriptsize, inner sep=2pt},                       % combinational logic
  hard/.style ={draw=harde, fill=hardf, thick, dashed, align=center,
                font=\scriptsize, inner sep=2pt},                       % hard macro / off-chip
  sram/.style ={draw=memd, fill=memf, thick, align=center, font=\scriptsize,
                inner sep=2pt},                                         % SRAM array, not flip-flops
  ff/.style   ={draw=#1, fill=#1!30, line width=0.35pt},                % one flip-flop square
  ffgone/.style={draw=muted, fill=white, densely dashed, line width=0.35pt},
  panel/.style={draw=muted, rounded corners=6pt, line width=0.9pt},
  ptitle/.style={anchor=north west, font=\normalsize\bfseries, text=ink, inner sep=0pt},
  note/.style ={font=\scriptsize, text=ink, align=left, inner sep=0pt},
  lab/.style  ={font=\tiny, text=muted, inner sep=1pt},
  axi/.style  ={{Stealth[length=2.6mm]}-{Stealth[length=2.6mm]}, line width=2.2pt, copper},
  axi1/.style ={-{Stealth[length=2.6mm]}, line width=2.2pt, copper},
  apb/.style  ={{Stealth[length=1.8mm]}-{Stealth[length=1.8mm]}, line width=1pt, copper},
  apb1/.style ={-{Stealth[length=1.8mm]}, line width=1pt, copper},
  dat/.style  ={-{Stealth[length=1.8mm]}, thick, pcbmid},
  ctl/.style  ={-{Stealth[length=1.8mm]}, densely dashed, thick, muted},
  irq/.style  ={-{Stealth[length=2mm]}, densely dotted, line width=1.3pt, bad},
  num/.style  ={circle, fill=copper, text=white, font=\scriptsize\bfseries, inner sep=1pt,
                minimum size=0.46cm},
  ours/.style ={draw=bad, line width=1.2pt, densely dashed, rounded corners=4pt},
}

% ---------------- helpers ----------------
\newcommand{\rn}[1]{\texttt{#1}}
% \ffrow{x}{y}{n}{cell}{style}: n flip-flops in a row, bottom-left corner at (x,y)
\newcommand{\ffrow}[5]{\foreach \q in {1,...,#3}{\draw[#5] ({#1+(\q-1)*#4},{#2}) rectangle ++(#4,#4);}}
% \ffcol{x}{y}{n}{cell}{style}: n flip-flops in a column, top-left corner at (x,y)
\newcommand{\ffcol}[5]{\foreach \q in {1,...,#3}{\draw[#5] ({#1},{#2-\q*#4}) rectangle ++(#4,#4);}}
% \ffgrid{x}{ytop}{n}{per row}{cell}{style}: n flip-flops, filled row by row from the top-left
\newcommand{\ffgrid}[6]{\pgfmathtruncatemacro{\gl}{#3-1}%
  \foreach \b in {0,...,\gl}{\pgfmathtruncatemacro{\gc}{mod(\b,#4)}\pgfmathtruncatemacro{\gr}{div(\b,#4)}%
    \draw[#6] ({#1+\gc*#5},{#2-(\gr+1)*#5}) rectangle ++(#5,#5);}}

% every figure: same page size, title, footer
\newcommand{\canvas}{\useasboundingbox (-0.3,-0.2) rectangle (38.6,30.4);}
\newcommand{\figtitle}[3]{%
  \node[num, minimum size=0.9cm, font=\Large\bfseries] at (0.45,29.55) {#1};
  \node[anchor=west, font=\LARGE\bfseries, text=pcb] at (1.15,29.55) {#2};
  \node[anchor=west, font=\small, text=ink] at (0,28.7) {#3};}
\newcommand{\figfoot}[1]{\node[anchor=west, font=\scriptsize, text=muted] at (0,0.25) {#1};}
% \panel{x0}{y0}{x1}{y1}{title}
\newcommand{\panel}[5]{\draw[panel] (#1,#2) rectangle (#3,#4);
  \node[ptitle] at ({#1+0.3},{#4-0.25}) {#5};}

% ---------------- timing diagrams: 1 cycle = \cw cm, low at y, high at y+0.45 ----------------
\newcommand{\cw}{1.0}
\newcommand{\wclk}[3]{% x0 y ncycles
  \foreach \k in {1,...,#3}{\draw[pcb, thick] ({#1+(\k-1)*\cw},#2) -- ++(0,0.45) -- ++(0.5*\cw,0)
     -- ++(0,-0.45) -- ++(0.5*\cw,0);}}
\newcommand{\wsig}[3]{% x0 y {level per cycle, e.g. 0,1,1,0}
  \gdef\wprev{x}%
  \foreach \v [count=\k from 0] in {#3}{%
    \draw[pcb, thick] ({#1+\k*\cw},{#2+0.45*\v}) -- ++(\cw,0);
    \if x\wprev\else\ifnum\wprev=\v\else\draw[pcb, thick] ({#1+\k*\cw},#2) -- ++(0,0.45);\fi\fi
    \global\let\wprev\v}}
\newcommand{\wbus}[5]{% x0 y c0 c1 text: one value held from cycle c0 to cycle c1
  \draw[pcbmid, thick, fill=white] ({#1+(#3)*\cw},{#2+0.225}) -- ({#1+(#3)*\cw+0.08},{#2+0.45})
    -- ({#1+(#4)*\cw-0.08},{#2+0.45}) -- ({#1+(#4)*\cw},{#2+0.225}) -- ({#1+(#4)*\cw-0.08},#2)
    -- ({#1+(#3)*\cw+0.08},#2) -- cycle;
  \node[font=\tiny\ttfamily] at ({#1+((#3)+(#4))*\cw/2},{#2+0.225}) {#5};}
\newcommand{\widle}[4]{% x0 y c0 c1: bus idle (no value)
  \draw[muted, thick] ({#1+(#3)*\cw},{#2+0.225}) -- ({#1+(#4)*\cw},{#2+0.225});}
\newcommand{\wname}[3]{\node[anchor=east, font=\scriptsize\ttfamily] at ({#1-0.15},{#2+0.225}) {#3};}

\begin{document}
\begin{tikzpicture}[line join=round]
\canvas
\figtitle{10}{Boot, security and debug: from power-on to the operating system}
  {A chip must start from code nobody can change, prove that every next piece of software is genuine,
   keep secrets away from normal apps, and still let engineers look inside.}

% =====================================================================
% (A) boot chain
% =====================================================================
\panel{0}{19.4}{38.3}{28.1}{Secure boot: a chain of trust}
\foreach \nm/\tx/\where/\sty [count=\k from 0] in {
    power-on/{PMIC raises the supply rails in order; reset is held until they are stable and the PLL is locked}/{--}/pwr,
    boot ROM/{CPU 0 starts at the reset vector in mask ROM (cannot be changed); reads boot pins + fuses}/{ROM}/sec,
    load BL1/{copy the first-stage loader from UFS / eMMC / SPI flash into on-chip SRAM (DRAM does not work yet)}/{ROM $\rightarrow$ SRAM}/sec,
    verify BL1/{hash + signature check against the key hash in the OTP fuses; bad $\Rightarrow$ stop or recovery mode}/{ROM}/sec,
    BL1 runs/{sets up clocks, trains the DRAM PHY (Figure 7): now DRAM can be used}/{SRAM}/mem,
    secure firmware/{load + verify the secure monitor (e.g. Trusted Firmware) and the TEE OS into secure DRAM}/{DRAM, secure}/sec,
    OS loader/{load + verify U-Boot / UEFI, then the Linux kernel + device tree}/{DRAM}/cpu,
    kernel + drivers/{drivers start; the NPU driver loads NPU firmware and requests its IRQ; apps start}/{DRAM}/cpu}{
  \pgfmathsetmacro{\x}{0.35+\k*4.7}
  \node[\sty, fill=white, minimum width=4.35cm, minimum height=4.6cm, text width=4.1cm, anchor=north west, font=\footnotesize, align=center]
    (b\k) at (\x,27.3) {\textbf{\k\ \nm}\\[4pt]\tx};
  \node[font=\tiny, text=muted, anchor=north] at ({\x+2.17},22.6) {runs from: \where};}
\foreach \k [evaluate=\k as \n using int(\k+1)] in {0,...,6}{\draw[dat] (b\k.east) -- (b\n.west);}
\draw[-{Stealth[length=2mm]}, thick, secd, bend right=6] ($(b1.south)+(0.6,-0.45)$) to ($(b4.south)+(-0.6,-0.45)$);
\foreach \k/\n in {4/5, 5/6, 6/7}{
  \draw[-{Stealth[length=2mm]}, thick, secd, bend right=15] ($(b\k.south)+(0.6,-0.45)$) to ($(b\n.south)+(-0.6,-0.45)$);}
\node[note, font=\small, anchor=north west, text width=37.6cm] at (0.35,21.2) {%
  \textcolor{secd}{\textbf{Red arrows}}: each stage checks the signature of the next one \textbf{before} jumping to it. The first link (the ROM and the key hash in the fuses)
  is fixed in silicon, so an attacker who changes the flash cannot run their own code. This is the \textbf{chain of trust}.};

% =====================================================================
% (B) signature check
% =====================================================================
\panel{0}{9.6}{12.6}{19.0}{How ``verify'' works}
\node[draw=pcb, fill=white, rounded corners=3pt, minimum width=3.6cm, minimum height=2.2cm, text width=3.4cm, align=center, font=\tiny, anchor=north west]
  (img) at (0.35,18.2) {\textbf{image in flash}\\[2pt]BL1 code\\+ public key\\+ signature};
\node[comb, minimum width=3.0cm, minimum height=1.0cm, text width=2.8cm, font=\tiny] (sha) at (6.0,17.6) {\textbf{SHA-256 engine}\\code $\rightarrow$ 256-bit digest};
\node[comb, minimum width=3.0cm, minimum height=1.0cm, text width=2.8cm, font=\tiny] (kh) at (6.0,16.1) {\textbf{hash of public key}\\$=$ value in OTP fuses?};
\node[comb, minimum width=3.0cm, minimum height=1.0cm, text width=2.8cm, font=\tiny] (sig) at (9.9,16.85) {\textbf{signature check}\\(ECDSA / RSA) with the key};
\draw[dat] (img.east |- sha) -- (sha); \draw[dat] (img.east |- kh) -- (kh);
\draw[dat] (sha.east) -- (sig.west |- sha.east); \draw[dat] (kh.east) -- (sig.west |- kh.east);
\node[good, draw=good, fill=good!15, rounded corners=3pt, font=\tiny\bfseries, text=good] at (8.3,14.4) {all match: jump to BL1};
\node[draw=bad, fill=bad!15, rounded corners=3pt, font=\tiny\bfseries, text=bad] at (11.1,14.4) {else: stop};
\draw[dat] (sig.south) -- (9.9,15.0) -| (8.3,14.65); \draw[dat] (9.9,15.0) -| (11.1,14.65);
\node[anchor=north west, font=\tiny\bfseries] at (0.35,13.6) {SHA-256 state = 8 registers $\times$ 32 bit:};
\ffgrid{0.4}{13.2}{256}{64}{0.14}{ff=secd}
\node[note, anchor=north west, text width=11.9cm] at (0.35,12.45) {%
  The hash engine is a datapath like your PE: 256 FF of state, plus adders, rotations and XORs between them, 64 rounds per 64-byte block.
  Changing one bit of the code changes the whole digest, so the signature no longer matches.\\[3pt]
  Only the private key (kept by the chip vendor, never on the chip) can make a valid signature.};

% =====================================================================
% (C) TrustZone
% =====================================================================
\panel{12.9}{9.6}{25.4}{19.0}{TrustZone: two worlds on one CPU}
\node[cpu, minimum width=5.5cm, minimum height=0.85cm, anchor=north west, font=\tiny] at (13.2,17.7) {EL0: apps};
\node[cpu, minimum width=5.5cm, minimum height=0.85cm, anchor=north west, font=\tiny] at (13.2,16.75) {EL1: Linux / Android kernel};
\node[cpu, minimum width=5.5cm, minimum height=0.85cm, anchor=north west, font=\tiny] at (13.2,15.8) {EL2: hypervisor (optional)};
\node[sec, minimum width=5.5cm, minimum height=0.85cm, anchor=north west, font=\tiny] at (19.5,17.7) {S-EL0: trusted apps (keys, payments, face data)};
\node[sec, minimum width=5.5cm, minimum height=0.85cm, anchor=north west, font=\tiny] at (19.5,16.75) {S-EL1: TEE OS};
\node[sec, minimum width=11.8cm, minimum height=0.85cm, anchor=north west, font=\tiny] at (13.2,14.8) {EL3: secure monitor: switches worlds when software executes an SMC instruction};
\node[font=\scriptsize\bfseries, text=cpud] at (15.95,18.0) {normal world};
\node[font=\scriptsize\bfseries, text=secd] at (22.25,18.0) {secure world};
\node[note, anchor=north west, text width=11.9cm] at (13.2,13.7) {%
  \textbf{Hardware does the separation}: every bus transfer carries an \textbf{NS bit} (non-secure; AxPROT[1] on AXI).\\[2pt]
  $\bullet$ \textbf{TZASC} sits in front of DRAM and blocks NS = 1 accesses to secure regions.\\
  $\bullet$ \textbf{TZPC} marks peripherals (crypto engine, fingerprint sensor) as secure-only.\\
  $\bullet$ Even a hacked Linux kernel cannot read secure memory: the bus returns an error.\\[3pt]
  An NPU can also be put in secure mode, e.g. to compare a face with a template kept in secure memory.};

% =====================================================================
% (D) security blocks
% =====================================================================
\panel{25.7}{9.6}{38.3}{19.0}{Security blocks on the die}
\foreach \nm/\tx [count=\k from 0] in {
    OTP fuses/{one-time programmable bits blown at the factory: root-key hash, device ID, ``debug off'', anti-rollback counter},
    crypto engine/{AES for bulk encryption, SHA for hashes, a big-number unit (PKA) for RSA / ECC},
    TRNG/{true random numbers: free-running ring oscillators jitter; the noise is sampled and cleaned by a hash},
    key store/{keys are used inside the crypto engine and never appear on the normal bus}}{
  \node[sec, fill=white, minimum width=12.0cm, minimum height=1.5cm, text width=11.7cm, anchor=north west, font=\scriptsize, align=left]
    at (26.0,{18.2-\k*1.75}) {\textbf{\nm}: \tx};}
\node[note, anchor=north west, text width=12.0cm] at (26.0,11.1) {%
  Attackers also measure power or inject glitches; secure blocks add noise, redundancy and sensors against that.};

% =====================================================================
% (E) debug
% =====================================================================
\panel{0}{0.7}{25.4}{9.3}{Debug: JTAG and the boundary-scan chain}
\node[per, fill=white, minimum width=3.0cm, minimum height=3.6cm, text width=2.8cm, anchor=north west, font=\tiny] (tap) at (0.35,8.5)
  {\textbf{TAP controller}\\16-state FSM\\driven by TMS\\[4pt]pins:\\TCK clock\\TMS mode\\TDI data in\\TDO data out};
\node[per, fill=white, minimum width=3.2cm, minimum height=0.8cm, font=\tiny, anchor=north west] at (3.7,8.5) {instruction register};
\node[per, fill=white, minimum width=3.2cm, minimum height=0.8cm, font=\tiny, anchor=north west] at (3.7,7.55) {IDCODE (32 bit)};
\node[per, fill=white, minimum width=3.2cm, minimum height=0.8cm, font=\tiny, anchor=north west] at (3.7,6.6) {BYPASS (1 bit)};
\node[per, fill=white, minimum width=3.2cm, minimum height=0.8cm, font=\tiny, anchor=north west] at (3.7,5.65) {debug port $\rightarrow$ CPU};
% chip with boundary scan ring
\draw[pcb, thick, rounded corners=3pt] (8.6,1.5) rectangle (15.6,8.5);
\node[font=\tiny, text=muted] at (12.1,5.0) {chip core};
\foreach \k in {0,...,9}{\draw[ff=prd] ({9.1+\k*0.6},7.9) rectangle ++(0.4,0.4);}
\foreach \k in {0,...,8}{\draw[ff=prd] (15.0,{7.3-\k*0.6}) rectangle ++(0.4,0.4);}
\foreach \k in {0,...,9}{\draw[ff=prd] ({14.5-\k*0.6},1.7) rectangle ++(0.4,0.4);}
\foreach \k in {0,...,8}{\draw[ff=prd] (8.8,{2.3+\k*0.6}) rectangle ++(0.4,0.4);}
\draw[dat] (7.4,8.1) -- (9.1,8.1);
\node[lab, text=ink, anchor=south] at (8.0,8.1) {TDI};
\draw[dat] (9.0,6.8) -- (7.4,6.8);
\node[lab, text=ink, anchor=south] at (8.0,6.8) {TDO};
\node[lab, text=ink, anchor=north, text width=6.6cm, align=center] at (12.1,1.4) {one scan flip-flop next to every pin, chained into a ring};
\node[note, anchor=north west, text width=9.3cm] at (15.9,8.5) {%
  \textbf{Boundary scan}: shift a pattern in on TDI, and the scan flip-flops drive the pins; capture what comes back and shift it
  out on TDO. That tests the solder joints on the board without running any software.\\[3pt]
  \textbf{Debug}: through the debug port a PC can halt a CPU core, single-step it, set breakpoints and read registers.
  Trace units record which instructions ran.\\[3pt]
  \textbf{In products}: a fuse turns debug off, or allows it only after a signed unlock.};

% =====================================================================
% (F) your FPGA
% =====================================================================
\panel{25.7}{0.7}{38.3}{9.3}{Your MAX 10: ``boot'' = configuration}
\node[note, anchor=north west, text width=12.0cm] at (26.0,8.5) {%
  An FPGA has no boot ROM program. At power-on it \textbf{loads its configuration}: the bits that set every LUT, flip-flop and wire.\\[4pt]
  $\bullet$ \textbf{.sof} through the USB-Blaster (JTAG, the same TAP as on the left) goes into configuration SRAM.
  Fast for testing; lost when the power goes off.\\[2pt]
  $\bullet$ \textbf{.pof} is written into the MAX 10's internal configuration flash. At every power-on the chip configures itself
  from it (``instant-on''); it can even hold two images.\\[4pt]
  $\bullet$ The same JTAG port is what SignalTap (Quartus' on-chip logic analyser) uses to watch your signals while the design runs.\\[4pt]
  Real SoCs add the security steps above; FPGAs can encrypt and authenticate the bitstream too (higher-end families).};

\figfoot{Boot stage names follow the common Arm flow (ROM, BL1, BL2 / TF-A, TEE, U-Boot, Linux). Exact stages differ between vendors.}
\end{tikzpicture}
\end{document}
